\usepackage{xcolor}
\definecolor{unipdred}{RGB}{155,0,30}
\definecolor{unipdblue}{RGB}{28,60,110}
\definecolor{unipdgray}{RGB}{95,95,95}
\definecolor{unipdgold}{RGB}{188,142,22}
\definecolor{unipdgreen}{RGB}{15,120,100} % teal-shifted green: stays distinguishable from unipdred under red-green color-vision deficiency


\usepackage[acronym,nowarn]{glossaries} % nowarn: this template prints acronyms via the separate `acronym` package instead, so \printglossary is deliberately never called
\makeglossaries          % ① only ONCE


%===================================
\usepackage[utf8]{inputenc}
\usepackage{minted}
\usepackage{fvextra}
\usepackage{upquote}
\usepackage{fancyvrb}
\setminted{escapeinside=||,mathescape=true}

\usepackage[backend=biber, style=apa, natbib=true, doi=true, url=false, isbn=false, eprint=false]{biblatex}


\addbibresource{references.bib}  % Replace with your own .bib filename

%===================================

\usepackage{newunicodechar}
\newunicodechar{Δ}{$\Delta$}
\newunicodechar{≈}{\approx}
\newunicodechar{−}{-}
\newunicodechar{⁻}{\textsuperscript{-}}
\newunicodechar{̃}{\~{}}
\usepackage{acronym}
\setlength{\headheight}{14.5pt}

%======================================================================
% --------------------- LANGUAGE AND ENCODING ------------------------
%======================================================================

\usepackage{textcomp} % For thin space and other text symbols
\DeclareUnicodeCharacter{2009}{\,} % Map Unicode thin space to LaTeX thin space
\DeclareUnicodeCharacter{03B1}{\alpha} % Map Unicode alpha to LaTeX \alpha
\usepackage[T1]{fontenc}             % Output font encoding for better character rendering
\usepackage[english]{babel}          % Language-specific typography and hyphenation
\usepackage[gen]{eurosym}            % Euro symbol support
\usepackage{amsthm} % For theorem, lemma, and remark environments
\usepackage{newunicodechar}

\DeclareUnicodeCharacter{202F}{~} % U+202F NARROW NO-BREAK SPACE as non-breaking space
\usepackage[most,breakable]{tcolorbox}
\usepackage{enumitem} % For custom enumerate labels
%======================================================================
% ------------------------- PAGE LAYOUT ------------------------------
%======================================================================
\usepackage[a4paper, margin=1in]{geometry} % Page size and margins
\usepackage{setspace}                      % For line spacing control (\onehalfspacing, \doublespacing)
\onehalfspacing                            % Set document to one-and-a-half line spacing

%======================================================================
% -------------------- FONTS AND TYPOGRAPHY --------------------------
%======================================================================
\usepackage{microtype}               % Improves typography (micro-adjustments to spacing)
\usepackage{csquotes}                % Context-sensitive quotation marks
\usepackage{empheq}                 % For empheq environment (boxed equations, etc.)

%======================================================================
% --------------------- CHAPTER AND SECTION STYLING ------------------
%======================================================================
\usepackage[Glenn]{fncychap}         % Stylized chapter headings
\ChNameVar{\bfseries\Large\sf\color{unipdred}}   % "Chapter" label in the institutional red
\ChNumVar{\Huge\color{unipdred}}                 % Chapter number
\ChTitleVar{\bfseries\Large\rm\color{unipdred}}  % Chapter title
\usepackage{titlesec}                % For customizing section/subsection titles
\titleformat{\section}{\large\bfseries}{\thesection}{1em}{}
\titleformat{\subsection}{\normalsize\bfseries}{\thesubsection}{1em}{}

%======================================================================
% -------------------- HEADER, FOOTER, AND FOOTNOTES -----------------
%======================================================================
\usepackage{fancyhdr}                % For custom headers and footers
\usepackage[bottom,hang]{footmisc}   % Advanced footnote placement
\pagestyle{fancy}
\fancyhf{}                           % Clear all header and footer fields
\renewcommand{\headrulewidth}{0.3pt} % Header rule line
\setlength{\footnotemargin}{6pt}     % Margin for footnote numbers

\renewcommand{\chaptermark}[1]{\markboth{\textit{\thechapter. #1}}{}}
\renewcommand{\sectionmark}[1]{\markright{\textit{\thesection. #1}}}

\fancyhead[L]{\ifnum\value{section}=0\leftmark\else\rightmark\fi} % Show Chapter or Section in left header
\fancyhead[R]{\thepage} % Page number in right header

% Style for chapter start pages (plain style)
\fancypagestyle{plain}{
    \fancyhf{} % Clear header and footer
    \renewcommand{\headrulewidth}{0pt} % No header line
    \fancyfoot[R]{\thepage} % Page number in bottom right
}


%======================================================================
% ------------------------- MATH PACKAGES ----------------------------
%======================================================================
\usepackage{amsmath, amssymb, amsthm} % Core AMS math packages and theorem environments

\usepackage{mathtools}               % Extends amsmath with more tools
\usepackage{bm}                      % For bold math symbols: \bm{}
\usepackage{siunitx}                 % For typesetting numbers and units correctly
\sisetup{output-decimal-marker={.}, per-mode=symbol}

%======================================================================
% ----------------------------- TABLES -------------------------------
%======================================================================
\usepackage{array}                   % Extended column types (m{}, b{}), used for vertically-centered wrapped columns
\usepackage{booktabs}                % High-quality table rules
\usepackage{multirow}                % For cells spanning multiple rows
\usepackage{tabularx}                % Tables with auto-adjusting column widths
\usepackage{threeparttable}          % For tables with notes
\usepackage{longtable}               % For tables that span multiple pages

%======================================================================
% ---------------------- FIGURES AND GRAPHICS ------------------------
%======================================================================
\usepackage{graphicx}                % For including images
\usepackage{caption}                 % Customizing figure and table captions
\captionsetup{labelfont={bf,color=unipdred}, font=small}
\usepackage{subcaption}              % For subfigures and sub-tables
\usepackage{float}                   % Improved control over float placement [H]

%======================================================================
% ------------------- DRAWING AND VISUALIZATION ----------------------
%======================================================================
\usepackage{tikz}                    % Powerful drawing language
\usepackage{pgfplots}                % For axis environment and advanced plots
\usetikzlibrary{shadows, arrows.meta, positioning, shapes.geometric}
\tikzset{li/.style={-latex, thick, dashed, color=gray}} % Define 'li' arrow style
\pgfplotsset{compat=1.18}
\usepgfplotslibrary{groupplots}
\usepgfplotslibrary{fillbetween}
\usepackage{pgfplotstable}           % For plotting data from tables/files
\usepackage{pgf-pie}                 % For creating pie charts
\usepackage{forest}                  % For drawing trees (e.g., decision trees)
\usepackage[table,x11names,svgnames]{xcolor}           % Color support, required by some packages

%======================================================================
% -------------------- CODE LISTINGS AND SYMBOLS ---------------------
%======================================================================
\usepackage{listings}                % For typesetting code
\lstset{basicstyle=\ttfamily, breaklines=true}
\usepackage{pifont}                  % Access to symbols like checkmarks (\ding{51})

%======================================================================
% -------------- HYPERLINKS AND CROSS-REFERENCING --------------------
%======================================================================
% NOTE: hyperref should be loaded last, with few exceptions.
\usepackage{hyperref}                % Creates hyperlinks in the document
\hypersetup{
    colorlinks=true,
    linkcolor=unipdblue,
    citecolor=unipdblue,
    urlcolor=unipdblue,
    pdfauthor={Qaiser Aziz},                                          % overridden again in main.tex from \ThesisAuthor
    pdftitle={Solar Grid Parity in Italy: Where Do We Stand}          % overridden again in main.tex from \ThesisTitleMain
}
\usepackage[capitalize,nameinlink]{cleveref} % Must be loaded AFTER hyperref

%======================================================================
% ----------------- THESIS STRUCTURE AND APPENDICES ------------------
%======================================================================
\usepackage[toc,page]{appendix}      % Better control over appendices
\usepackage[nottoc]{tocbibind}       % Adds Bib, LoF, LoT, etc., to the ToC

%======================================================================
% ----------------------- OPTIONAL UTILITIES -------------------------
%======================================================================
\usepackage{pdfpages}                % To insert pages from external PDF files
\usepackage{comment}                 % For commenting out large blocks of text
\usepackage{lipsum}                  % For generating dummy text (for testing layout)

%======================================================================
% ------------------------- END OF PREAMBLE --------------------------
%======================================================================
